% =====================================================================
%  FOUNDATIONS: MATHEMATICAL TOOLKIT  (precedes every volume)
% =====================================================================
\foundationschapter{Mathematical Toolkit}\chlabel{ch:toolkit}

Physics is stated in mathematics, and a derivation that keeps stopping to justify its algebra is
hard to follow. This chapter collects, in one place, the mathematical tools the rest of the book
relies on, so that later arguments need not break off to explain a technique. Nothing here is examined
for its own sake. What is examined, constantly, is whether you can reach for the right tool at the
right moment, and that is a skill best built before the physics arrives to distract you.

There are sixteen tools, and each one earns its place by appearing in a later chapter; each entry ends
with a pointer to where it is used. The first twelve serve mechanics (Volume~I) and everything after it;
the last four, exponentials and logarithms, first-order rate equations, averages over a distribution and
probability per unit time, are first needed for heat and gases (Volume~II) and for radioactivity
(Volume~III), and can be left until then. Read the chapter once now, quickly, for orientation. Then return
to it whenever a step in a derivation feels unfamiliar. The test of whether you have understood a
tool is not whether you can state it but whether you can say what it is \emph{for}: the dot product
answers ``how much along?'', the derivative answers ``how fast, right now?'', the integral answers
``how much, in total?''.

One honest warning. The calculus here is used, not proved. Where a result such as the derivative of
$\sin x$ is quoted, a mathematics course will supply the proof; this book supplies the physical
meaning and the habit of use. That division of labour is deliberate, and it is the standard one.

\section{Angles, Radians and the Trigonometric Ratios}

Two facts about angles are used on almost every page of this book: that calculus requires them to be
measured in radians, and that trigonometry is mostly used to split one vector into two perpendicular parts.
(A \emph{vector} is a quantity with a direction as well as a size, such as a displacement or a force;
its size is its \emph{magnitude}. \chref{ch:vectors} makes the idea precise.)

\begin{defbox}{Radians: the only angle measure that works in calculus}
An angle in \textbf{radians} is defined as the arc it subtends divided by the radius:
\[
\theta=\frac{\text{arc}}{\text{radius}}=\frac{s}{R}\qquad\Longrightarrow\qquad s=R\theta.
\]
Because it is a length divided by a length, a radian is \emph{dimensionless}. What makes it the only
measure fit for calculus is that, in radians and only in radians, $\sin\theta\approx\theta$ for small
$\theta$ and $\frac{d}{d\theta}\sin\theta=\cos\theta$; in degrees every such result carries an extra
factor $\pi/180$. A full turn is $2\pi$ radians, so
\[
180^\circ=\pi\ \text{rad},\qquad 1\ \text{rad}\approx57.3^\circ.
\]
\emph{Used in:} arc length and angular displacement (\chref{ch:circkin}), $v=\omega r$ (\chref{ch:circkin}), and every
rotational formula in \chsref{ch:rotation}{ch:rolling}. \textbf{An angle that stands on its own in a
formula} (as in $s=R\theta$ or $\omega=\theta/t$), rather than inside $\sin$, $\cos$ or $\tan$,
\textbf{must be in radians}; convert first. Inside a trigonometric function either unit is safe,
provided the calculator is set to match.
\end{defbox}

\begin{defbox}{The trigonometric ratios, and resolving a vector}
In a right triangle with hypotenuse of length $A$ and an angle $\theta$ (the sides named relative to $\theta$),
\[
\sin\theta=\frac{\text{opposite}}{\text{hyp}},\qquad
\cos\theta=\frac{\text{adjacent}}{\text{hyp}},\qquad
\tan\theta=\frac{\sin\theta}{\cos\theta},
\]
together with the identity that follows from Pythagoras,
\[
\boxed{\;\sin^{2}\theta+\cos^{2}\theta=1\;}
\]
\textbf{The one use that dominates this book} is \emph{resolution}: replacing a vector by two
perpendicular pieces that can be handled independently, the two legs of a right triangle whose hypotenuse
is the vector (Figure~\ref{fig:tkcomp}).
\[
A_x=A\cos\theta,\qquad A_y=A\sin\theta,\qquad
A=\sqrt{A_x^{2}+A_y^{2}},\qquad \tan\theta=\frac{A_y}{A_x},
\]
where $\theta$ is measured from the $x$-axis and $A_x$, $A_y$ are the \emph{components} along the $x$- and
$y$-axes. \emph{Used in:} projectile motion (\chref{ch:projectiles}), every inclined-plane problem
(Chapters~\ref{ch:newton}--\ref{ch:friction}), and the components
of weight on a slope, $mg\sin\theta$ and $mg\cos\theta$.
\end{defbox}

\begin{figure}[htbp]
\centering
\tikzfig{tk_components}
\figcaption{fig:tkcomp}{Resolution is just a right triangle read backwards: the vector is the
hypotenuse, and its components are the two legs. Note which leg gets the cosine; it is always the one
\emph{adjacent} to the marked angle, so the labelling depends on where $\theta$ is measured from.}
\end{figure}

\begin{trapbox}{$\cos$ does not always go with the horizontal}
$A_x=A\cos\theta$ only when $\theta$ is measured \emph{from the $x$-axis}. On a plane inclined at $\theta$
to the horizontal, the weight makes the angle $\theta$ with the \emph{normal} to the slope, not with the
slope itself, so the component along the slope is the one opposite $\theta$: $mg\sin\theta$. Always ask
``adjacent or opposite to the angle I have marked?'' rather than reciting a formula.
\end{trapbox}

\section{Vector Algebra}\label{sec:vecalg}

Adding vectors is the business of \chref{ch:vectors}. What needs flagging now is that vectors (written
$\vc{A}$, with magnitude $A$) can be multiplied in two genuinely different ways, answering two different
questions. Knowing which question a situation asks is more useful than knowing either formula.

\begin{defbox}{The two products, and what each is for}
\textbf{Scalar (dot) product:} returns a \emph{number}:
\[
\vc{A}\cdot\vc{B}=AB\cos\theta=A_xB_x+A_yB_y+A_zB_z,
\]
where $\theta$ is the angle between the vectors drawn tail to tail and $A_x$, $A_y$, $A_z$ are the
components along three perpendicular axes. It measures how much of one vector lies \emph{along} the
other, and for non-zero vectors it vanishes exactly when they are perpendicular. \emph{Used in:} work $W=\vF\cdot\vc{d}$ and power $P=\vF\cdot\vv$ (\chref{ch:energy}).

\medskip
\textbf{Vector (cross) product:} returns a \emph{vector} $\vc{A}\ex\vc{B}$ of magnitude $AB\sin\theta$,
perpendicular to both $\vc{A}$ and $\vc{B}$, its sense fixed by a right-hand rule. Its magnitude is the
\emph{area} of the parallelogram they span, and it vanishes when the vectors are parallel. (The
component forms, the right-hand rule and worked uses of both products are in \S\ref{sec:vecprod}.) \emph{Used in:} torque $\vtau=\vr\ex\vF$ and angular
momentum $\vL=\vr\ex\vp$ (\chref{ch:rotation}).

\medskip
\textbf{The memory hook:} the dot product asks ``how much \emph{along}?'', the cross product ``how much
\emph{across}?''
\end{defbox}

\section{Differentiation}

Most quantities in mechanics are rates: velocity is the rate of change of position, power the rate of doing
work, force the rate of change of momentum. The derivative is the mathematical form of a rate, and only a
handful of its rules are needed.

\begin{defbox}{The derivative is a rate, and a slope}
The derivative of $y$ with respect to $x$ is the limit of a ratio of small changes, where $\Delta x$
(read ``delta $x$'') is a change in $x$ and $\Delta y$ the change in $y$ it produces:
\[
\frac{dy}{dx}=\lim_{\Delta x\to0}\frac{\Delta y}{\Delta x}.
\]
Geometrically it is the \textbf{slope of the tangent} to the graph, the limiting position of a secant
whose two points close in on each other (Figure~\ref{fig:tkderiv}). Physically it is an
\textbf{instantaneous rate}: how fast $y$ changes right now, as opposed to the average change over an
interval. Its units are those of $y$ divided by those of $x$, which is why a rate of change of position
is measured in \si{m/s}. For example, if $x=5t^{2}$ (metres, seconds) then $dx/dt=10t$, so at
$t=\SI{3}{s}$ the instantaneous rate is \SI{30}{m/s}, while the average over the first three seconds is
only $\SI{45}{m}/\SI{3}{s}=\SI{15}{m/s}$. Differentiating twice gives the \textbf{second derivative}
$d^{2}y/dx^{2}$, the rate of change of the rate; acceleration, $d^{2}x/dt^{2}$, is the first example.

The derivatives this book needs:
\[
\frac{d}{dx}\,x^{n}=nx^{n-1},\qquad
\frac{d}{dt}\sin\omega t=\omega\cos\omega t,\qquad
\frac{d}{dt}\cos\omega t=-\omega\sin\omega t,
\]
\[
\frac{d}{dx}e^{kx}=k\,e^{kx},\qquad \frac{d}{dx}\ln x=\frac{1}{x}.
\]
\emph{Used in:} velocity $\vv=d\vr/dt$ and acceleration $\va=d\vv/dt$ (\chref{ch:straightline}), instantaneous
power $P=dW/dt$ (\chref{ch:energy}), and $\vF=d\vp/dt$ (\chref{ch:newton}).
\end{defbox}

\begin{figure}[htbp]
\centering
\tikzfig{tk_derivative}
\figcaption{fig:tkderiv}{Average versus instantaneous. The green \emph{secant} joins two separated points
and gives the average rate $\Delta s/\Delta t$. Sliding the second point inward until it merges with the
first turns the secant into the red \emph{tangent}, whose slope is $ds/dt$, the instantaneous rate.
That limiting process is the whole content of the derivative.}
\end{figure}

\begin{methodbox}{The two combination rules we need}
\textbf{Product rule:}
\[
\frac{d}{dx}(uv)=u\frac{dv}{dx}+v\frac{du}{dx}
\]
\textbf{Chain rule} (for a function of a function: differentiate the outside, then multiply by the
derivative of the inside; a prime denotes a derivative, $f'=df/du$):
\[
\frac{d}{dx}f\big(u(x)\big)=f'\big(u(x)\big)\cdot u'(x)
\qquad\text{for example}\qquad
\frac{d}{dt}\sin(\omega t)=\omega\cos(\omega t).
\]
\emph{Used in:} the product rule for differentiating $\vv=v\,\uhat{t}$, speed times a unit vector, which
splits acceleration into its tangential and centripetal parts (\chref{ch:circkin}); the chain rule for
every sinusoid whose argument is not just $t$, and for $a=dv/dt=v\,dv/dx$ (\chref{ch:straightline}).
\end{methodbox}

\begin{defbox}{Maxima and minima: where the slope is zero}
At the top of a hill or the bottom of a valley the tangent is horizontal. So a smoothly varying
quantity $y(x)$ is largest or smallest where
\[
\frac{dy}{dx}=0.
\]
To tell which, look at how the slope changes: if $d^{2}y/dx^{2}<0$ the slope is falling from positive to
negative, and the point is a \textbf{maximum}; if $d^{2}y/dx^{2}>0$ it is a \textbf{minimum}. For
$y=4x-x^{2}$, $dy/dx=4-2x$ vanishes at $x=2$, and $d^{2}y/dx^{2}=-2<0$, so the greatest value is
$y(2)=4$. Two cautions: a zero slope can also be a mere pause ($y=x^{3}$ at $x=0$), and when $x$ is
confined to a range the largest value may sit at an end of it, where the slope need not vanish.
\emph{Used in:} the greatest velocity under a varying acceleration (\chref{ch:straightline}), the least
drift in crossing a river (\chref{ch:relative}), the least force to drag a block (\chref{ch:friction}),
and stable and unstable equilibrium (\chref{ch:energy}).
\end{defbox}

\begin{methodbox}{Logarithmic differentiation: where the error rules come from}
Taking a logarithm turns products into sums, and \emph{then} differentiating gives a rule for
\emph{relative} changes. If
\[
Z=\frac{A^{p}B^{q}}{C^{r}},\qquad\text{then}\qquad \ln Z=p\ln A+q\ln B-r\ln C,
\]
and differentiating each term (using $d(\ln x)=dx/x$, where $dA$ means a very small change in $A$)
gives at once
\[
\frac{dZ}{Z}=p\frac{dA}{A}+q\frac{dB}{B}-r\frac{dC}{C}.
\]
For small errors the differentials can be replaced by the errors themselves, $dA\to\Delta A$ and so on.
The errors in $A$, $B$ and $C$ may have either sign, so in the worst case all three terms push the same
way and their magnitudes add (with $p,q,r$ positive):
\[
\boxed{\;\frac{\Delta Z}{Z}=p\frac{\Delta A}{A}+q\frac{\Delta B}{B}+r\frac{\Delta C}{C}\;}
\]
\emph{Used in:} the combination of errors (\chref{ch:measurement}); this single line \emph{is} the derivation of the
rule that percentage errors add, with powers as multipliers.
\end{methodbox}

\section{Integration}\label{sec:integration}

Integration is used in this book almost entirely in one way: to add up the contributions of many small
pieces of a body, each too small to matter alone. The picture to hold is a sum of thin slices.

\begin{defbox}{An integral is a sum, and an area}
Integration undoes differentiation, but the picture that matters here is different: an integral is the
limit of a \textbf{sum of infinitely many thin slices}. Cutting a region into strips of width $\Delta x$
and height $f(x)$ and adding their areas gives $\sum f(x)\Delta x$; letting the strips become
infinitesimally thin (Figure~\ref{fig:tkint}) turns the sum between $x=a$ and $x=b$ into the \textbf{definite integral}
\[
\int_a^b f(x)\,dx=\lim_{\Delta x\to0}\sum f(x)\,\Delta x,
\]
the area under the graph. The link to differentiation, quoted here from mathematics, is what makes such
sums computable: if $G(x)$ is any function whose derivative is $f(x)$, then
$\int_a^b f(x)\,dx=G(b)-G(a)$. Without limits, the \textbf{indefinite integral} $\int f\,dx=G(x)+C$
carries an arbitrary constant $C$ (not the $C$ of the error rule above), since adding a constant to $G$ leaves its derivative unchanged; in
physics $C$ is fixed by a known starting value, such as the position at $t=0$. For example, $\int_0^{\pi}\sin\theta\,d\theta=[-\cos\theta]_0^{\pi}=2$.
The integrals this book needs:
\[
\int x^{n}\,dx=\frac{x^{n+1}}{n+1}\ (n\neq-1),\qquad
\int_0^{\pi}\sin\theta\,d\theta=2,\qquad
\int_0^{\pi}\sin^{3}\theta\,d\theta=\frac{4}{3}.
\]
\emph{Used in:} work by a variable force $W=\int F\,dx$ (\chref{ch:energy}), and every centre-of-mass and
moment-of-inertia calculation, where the ``slices'' are literal pieces of the body
(\chsref{ch:com}{ch:rotation}).
\end{defbox}

\begin{figure}[htbp]
\centering
\tikzfig{tk_integral}
\figcaption{fig:tkint}{Why an integral is an area. Coarse strips (left) already approximate the region;
making them finer (right) makes the approximation exact in the limit. This is the picture to hold when a
derivation says ``take an element $dm$ and integrate''; you are adding up slices.}
\end{figure}

\begin{methodbox}{The pattern behind every $\int r^2\,dm$}
Whenever this book integrates over a body, three steps recur:
\begin{enumerate}
\item \textbf{Choose a slice} whose every point shares the same value of the quantity being weighted
      (the same distance $r$ from an axis, or the same height $y$).
\item \textbf{Express its mass} $dm$ in terms of a density and the slice's size: $\lambda\,dx$ for a
      wire ($\lambda$ the mass per unit length), $\sigma\,dA$ for a lamina ($\sigma$ the mass per unit
      area, $dA$ here the slice's area), $\rho\,dV$ for a solid ($\rho$ the mass per unit volume).
\item \textbf{Integrate} over the range that sweeps the slice through the whole body.
\end{enumerate}
Choosing the slice well is the entire art; the integration itself is then routine.
\end{methodbox}

\section{Approximations and Algebra}\label{sec:approx}

The next group of tools turns equations into numbers. None is deep, but each comes with a habit of
interpretation that matters as much as the algebra.

\begin{defbox}{The small-angle approximation}
For a small angle measured \textbf{in radians},
\[
\boxed{\;\sin\theta\approx\theta,\qquad \tan\theta\approx\theta,\qquad \cos\theta\approx1-\tfrac12\theta^{2}\approx1\;}
\]
The reason is visible geometrically: for a thin wedge the arc $R\theta$, the opposite side $R\sin\theta$
and the tangent $R\tan\theta$ all shrink towards one another (Figure~\ref{fig:tksmall}, left). The error in $\sin\theta\approx\theta$
stays below one percent up to about $14^\circ$, and in $\tan\theta\approx\theta$ up to about $10^\circ$;
the cruder $\cos\theta\approx1$ is already one percent out at $8^\circ$.
\emph{Used in:} the derivative of a rotating unit vector (\chref{ch:circkin}), the small-angle limit of
the conical pendulum (\chref{ch:circdyn}), and anywhere a curved arc is replaced by a straight chord.
\end{defbox}

\begin{figure}[htbp]
\centering
\tikzfig{tk_smallangle}
\figcaption{fig:tksmall}{Left: in a wedge of angle $\theta$ the perpendicular $R\sin\theta$, the arc $R\theta$ and the tangent $R\tan\theta$ always come in that order, and they close up on one another as the wedge narrows. Right: the three graphs leave the origin with the same slope, which is exactly why the approximation works. In the shaded band ($\theta<10^\circ$) all three agree within one percent; beyond it they part, $\tan\theta$ fastest.}
\end{figure}

\begin{defbox}{The binomial approximation}
For any power $n$ (whole, negative or fractional) and a small number $x$, $|x|\ll1$,
\[
\boxed{\;(1+x)^{n}\approx1+nx\;}
\]
For a whole number $n$ this is the start of the binomial expansion,
$(1+x)^{n}=1+nx+\tfrac12n(n-1)x^{2}+\cdots$, in which every later term carries a further factor of the
small $x$ and can be dropped; that the same first two terms hold for any $n$ is quoted from mathematics.
Geometrically it replaces the curve $(1+x)^{n}$ by its tangent at $x=0$, whose slope is $n$. For example,
$\sqrt{1.02}=(1+0.02)^{1/2}\approx1.01$ (true value $1.00995$), and $1/1.01^{2}\approx1-0.02=0.98$ (true
value $0.9803$). It is the same idea as the error rule above: $(1+\Delta A/A)^{p}\approx1+p\,\Delta A/A$.
\emph{Used in:} the fall of $g$ with height, $(1+h/R)^{-2}\approx1-2h/R$ (\chref{ch:gravitation}),
valid only while $h\ll R$.
\end{defbox}

\begin{keybox}{The quadratic formula}
\[
ax^{2}+bx+c=0\qquad\Longrightarrow\qquad x=\frac{-b\pm\sqrt{b^{2}-4ac}}{2a}
\]
When a physics problem yields two roots, \textbf{both must be examined}: often both are meaningful
(the launch and the landing), and sometimes one is negative or unphysical and is rejected for a stated
reason. Deciding which to keep is part of the physics. \emph{Used in:} finding the time of flight from
$y=ut-\tfrac12gt^{2}$ (\chref{ch:projectiles}), where one root is the launch instant and the other the landing.
\end{keybox}

\begin{methodbox}{Simultaneous equations, and eliminating the unknown you do not want}
Connected-body problems produce one equation per body sharing a common unknown. The usual technique:
\textbf{add or subtract the equations so that the unwanted unknown cancels.} For the Atwood machine
(two masses hung over a pulley, $m_1>m_2$, joined by a string of tension $T$),
\[
m_1g-T=m_1a,\qquad T-m_2g=m_2a,
\]
and \emph{adding} them eliminates $T$ in one step, leaving $(m_1-m_2)g=(m_1+m_2)a$. \emph{Used in:} pulleys
(\chref{ch:connected}) and elastic collisions (\chref{ch:collisions}), where dividing the energy equation by the momentum equation
eliminates the messy squares.
\end{methodbox}

\begin{methodbox}{Proportional reasoning: the fastest tool of all}
Many exam questions ask only \emph{how a quantity scales}, not its value. If $y\propto x^{n}$ then
multiplying $x$ by $k$ multiplies $y$ by $k^{n}$:
\begin{itemize}
\item $P\propto v^{3}$ (the power to push through air whose drag force grows as $v^2$, since power is
      force times speed): double the speed, \emph{eight times} the power.
\item $K=\tfrac12mv^{2}$: double the speed, \emph{four times} the kinetic energy, and, for the same
      braking force, four times the braking distance.
\item $I\propto MR^{2}$ at fixed mass: double the radius, \emph{four times} the moment of inertia.
\end{itemize}
No numbers are needed, only the exponent. \emph{Used in:} drag and power (\chref{ch:energy}), stopping
distances (\chref{ch:straightline}), moments of inertia (\chref{ch:rotation}), and Kepler's third law
(\chref{ch:kepler}).
\end{methodbox}

\section{Exponentials and Logarithms}\label{sec:explog}

Mechanics rarely needs the exponential, but the rest of physics needs it constantly, and always for the
same reason. Whenever a quantity changes at a rate proportional to its own size, the exponential is the
answer: a hot body cools faster the hotter it is, a radioactive sample loses more nuclei per second the
more nuclei it has, and a population of neutrons in a reactor grows faster the larger it already is.
This section builds the exponential from that one property, introduces the logarithm as the tool for
undoing it, and then solves the two rate equations that the later volumes keep meeting.

\begin{defbox}{The exponential: the function equal to its own rate}
The \textbf{exponential function} $e^{x}$ is defined by two properties: it equals its own derivative,
and it is $1$ at $x=0$,
\[
\frac{d}{dx}e^{x}=e^{x},\qquad e^{0}=1.
\]
The number $e=2.71828\ldots$ is its value at $x=1$. It can be met without calculus as the limit of
compound growth: a sum growing by $100\%$ in a year, but credited in $n$ equal instalments of
$(100/n)\%$, becomes $(1+1/n)^{n}$ times larger, which is $2.613$ for monthly credit, $2.715$ for daily,
and tends to $e$ as $n\to\infty$. The rules of powers carry over, because $e^{x}$ \emph{is} a power of
$e$:
\[
e^{a+b}=e^{a}e^{b},\qquad e^{-a}=\frac{1}{e^{a}},\qquad (e^{a})^{b}=e^{ab}.
\]
By the chain rule, $\dfrac{d}{dt}e^{kt}=k\,e^{kt}$ for any constant $k$: the rate of $e^{kt}$ is $k$ times
itself. So $e^{kt}$ grows for $k>0$ and decays towards zero for $k<0$, never reaching it. For small $x$
the curve is close to its tangent at the origin, whose slope is $e^{0}=1$, so
\[
e^{x}\approx1+x\qquad(|x|\ll1),
\]
the exponential relative of the binomial approximation of Section~\ref{sec:approx}; for example
$e^{-0.1}=0.905$ against $1-0.1=0.900$.
\end{defbox}

\begin{defbox}{The natural logarithm}
The \textbf{natural logarithm} $\ln x$ (defined for $x>0$) undoes the exponential: $y=\ln x$ means
$x=e^{y}$. So $\ln1=0$, $\ln e=1$, and $\ln x$ is negative for $0<x<1$. Because exponents add when
powers multiply, logarithms turn products into sums and powers into multipliers:
\[
\ln(ab)=\ln a+\ln b,\qquad \ln\frac{a}{b}=\ln a-\ln b,\qquad \ln a^{p}=p\ln a.
\]
Two values recur so often that they are worth memorising: $\ln2=0.693$ and $\ln10=2.303$. The
logarithm to base $10$, written $\log_{10}$, is related by $\log_{10}x=\ln x/\ln10$; calculators label
the natural one ``ln'' and the base-$10$ one ``log''.

The derivative of $\ln x$, quoted in the differentiation toolkit above, follows from that of $e^{y}$. If
$y=\ln x$ then $x=e^{y}$, so $dx/dy=e^{y}=x$, and turning the rate upside down gives
\[
\frac{d}{dx}\ln x=\frac{1}{x}\qquad\Longrightarrow\qquad \int_a^b\frac{dx}{x}=\ln b-\ln a=\ln\frac{b}{a}.
\]
That integral is the one gap in the power rule $\int x^{n}dx=x^{n+1}/(n+1)$, which fails at $n=-1$:
the area under $1/x$ is a logarithm, and it depends only on the \emph{ratio} $b/a$.
\end{defbox}

\begin{derivbox}{First-order decay: $dy/dt=-ky$}
\textbf{The equation.} Suppose a quantity $y(t)$ decreases at a rate proportional to itself, with a
constant $k>0$ (units \si{s^{-1}}):
\[
\frac{dy}{dt}=-ky .
\]
Equivalently, the \emph{fractional} rate of loss, $-(dy/dt)/y$, is the same constant $k$ whatever the
size of $y$. This is the property to recognise: a fixed fraction lost per second.

\stepn{1} \textbf{Separate the variables.} Divide by $y$ and multiply by $dt$, so that each side holds
one variable: $dy/y=-k\,dt$.

\stepn{2} \textbf{Integrate} from the starting value $y_0$ at $t=0$ to the value $y$ at time $t$, using
$\int dy/y=\ln y$:
\[
\ln\frac{y}{y_0}=-kt\qquad\Longrightarrow\qquad \boxed{\;y=y_0\,e^{-kt}\;}
\]
\stepn{3} \textbf{Check.} Differentiating, $dy/dt=-k\,y_0e^{-kt}=-ky$, and at $t=0$, $y=y_0$. Both the
equation and the starting value are satisfied, so this is the solution.

\stepn{4} \textbf{Two natural clocks.} The \textbf{time constant} $\tau=1/k$ is the time in which $y$
falls by the factor $e$, to $0.368\,y_0$. The \textbf{half-life} $t_{1/2}$ is the time in which it falls
to half; setting $e^{-kt_{1/2}}=\tfrac12$ gives
\[
t_{1/2}=\frac{\ln2}{k}=0.693\,\tau .
\]
Since $e^{-kt}=\big(e^{-kt_{1/2}}\big)^{t/t_{1/2}}$, the same law reads $y=y_0\,2^{-t/t_{1/2}}$: after $n$
half-lives the fraction $(\tfrac12)^{n}$ remains, whatever the starting value.
\end{derivbox}

The same equation with $+k$ in place of $-k$ describes unchecked growth, $y=y_0e^{kt}$, which doubles in
every interval $\ln2/k$. Real growth is usually held in check by a loss that grows with $y$, and that
leads to the second equation.

\begin{derivbox}{Growth with supply: $dy/dt=P-ky$}
\textbf{The equation.} Now let $y$ be supplied at a constant rate $P$ (units of $y$ per second) while it is
lost at the rate $ky$:
\[
\frac{dy}{dt}=P-ky .
\]
\stepn{1} \textbf{Find the balance.} The rate is zero, and $y$ stops changing, when supply equals loss:
$y_\infty=P/k$. This \emph{equilibrium} (or saturation) value is where $y$ heads.

\stepn{2} \textbf{Measure from the balance.} Let $w=y-P/k$, the distance still to go. Since $P/k$ is
constant, $dw/dt=dy/dt=P-ky=-k\,w$: the gap obeys plain first-order decay, so $w=w_0e^{-kt}$.

\stepn{3} \textbf{Return to $y$.} Adding $P/k$ back,
\[
\boxed{\;y=\frac{P}{k}+\Big(y_0-\frac{P}{k}\Big)e^{-kt}\;}
\qquad\text{and, from }y_0=0,\qquad y=\frac{P}{k}\big(1-e^{-kt}\big).
\]
\stepn{4} \textbf{Read it.} The gap to saturation halves in every half-life $\ln2/k$, so starting from
zero $y$ reaches $\tfrac12$, $\tfrac34$, $\tfrac78$ of $P/k$ after one, two and three half-lives. The time
constant is the same $\tau=1/k$ as in pure decay: the loss mechanism alone sets the clock, while the
supply sets only the level. Supply for many half-lives gains almost nothing; the level is already
saturated.
\end{derivbox}

\begin{figure}[htbp]
\centering
\tikzfig{tk_exponential}
\figcaption{fig:tkexp}{Left: first-order decay (blue) halves in every half-life, so the marked values are
$\tfrac12$, $\tfrac14$, $\tfrac18$; growth with supply (teal) climbs towards its saturation level $P/k$,
closing half of the remaining gap in every half-life. The two curves are mirror images about the
half-way line, because the gap in the second obeys the first. Right: on a logarithmic vertical axis the
decay is a straight line of slope $-k$, the signature by which an exponential is recognised in data.}
\end{figure}

\begin{defbox}{Semilog plots: how an exponential is recognised}
Taking the logarithm of $y=y_0e^{-kt}$ gives
\[
\ln y=\ln y_0-kt,
\]
a \emph{straight line} when $\ln y$ (or $y$ on a logarithmic scale) is plotted against $t$
(Figure~\ref{fig:tkexp}, right). Its slope is $-k$, so two points read off the line give
\[
k=\frac{\ln(y_1/y_2)}{t_2-t_1}.
\]
Plotting this way answers two questions at a glance: whether a set of measurements is exponential at all
(the points lie on a line), and with what rate (the slope). A curved semilog plot is informative too: a
sum of two exponentials with different $k$ is not a straight line, steep at first while the faster one
dies away, then straight with the slower one's slope.
\end{defbox}

\begin{exbox}{Using the decay law with logarithms}
A quantity decays as $y=y_0e^{-kt}$ with half-life $t_{1/2}=\SI{8.0}{h}$. (a) Find $k$ and the time
constant. (b) What fraction remains after one day? (c) How long until only $1\%$ remains?

\stepn{1} $k=\ln2/t_{1/2}=0.693/\SI{8.0}{h}=\SI{0.0866}{h^{-1}}$, and $\tau=1/k=\SI{11.5}{h}$, longer than
the half-life, as it must be.

\stepn{2} One day is $24/8=3$ half-lives, so the fraction left is $(\tfrac12)^{3}=\tfrac18=0.125$. No
logarithm is needed when the time is a whole number of half-lives.

\stepn{3} For $1\%$ the number of half-lives is not whole, so take logarithms: $e^{-kt}=0.01$ gives
$t=\ln100/k=4.605/0.0866=\SI{53}{h}$, about $6.6$ half-lives.
\ans{\SI{0.0866}{h^{-1}}, \SI{11.5}{h}; $0.125$; \SI{53}{h}}
\end{exbox}

\begin{exbox}{Growth with supply: how long to saturate?}
A quantity is supplied at $P=500$ units per second and lost at the rate $ky$ with $k=\SI{0.020}{s^{-1}}$,
starting from zero. Find the saturation level, the level after \SI{100}{s}, and the time to reach $90\%$
of saturation.

\stepn{1} Saturation: $P/k=500/0.020=25\,000$ units.

\stepn{2} After \SI{100}{s}, $kt=2$, so $y=25\,000\,(1-e^{-2})=25\,000\times0.865=21\,600$ units.

\stepn{3} $90\%$ means a remaining gap of $10\%$: $e^{-kt}=0.1$, so $t=\ln10/k=2.303/0.020=\SI{115}{s}$.
\ans{$25\,000$; $21\,600$; \SI{115}{s}}
\end{exbox}

\begin{trapbox}{an exponential never ``finishes''}
After two half-lives a quarter remains, not none; after ten, about $0.1\%$ still remains. A half-life is
not half of a lifetime. Two further slips are common: confusing the half-life $t_{1/2}=0.693/k$ with the
time constant $\tau=1/k$ (they differ by the factor $\ln2$), and taking a logarithm with the
``log'' key, which gives base $10$, where the natural logarithm is meant.
\end{trapbox}

\textbf{Where the model stops working.} The solutions above assume that $k$ (and $P$) stay constant. A
cooling rate constant that itself changes with temperature, or a supply that is switched off, needs a new
calculation for each stage. And when $y$ counts discrete things, such as a handful of nuclei, the smooth
curve describes only the \emph{average} behaviour; the individual events are random, which is the
subject of the next section.

\medskip\noindent
\emph{Used in:} Newton's law of cooling, where the temperature excess decays exponentially
(\chref{ch:thermal}, where the rate constant is written $\kappa$); isothermal work, where integrating $dV/V$ gives the logarithm of the volume ratio (\chref{ch:thermo});
the radioactive decay law, half-life and mean life, and production at a constant rate
$dN/dt=\Prod-\lambda N$ (\chref{ch:decaylaw}); the exponential growth of a chain reaction
(\chref{ch:reactions}); and radioactive dating and the absorption of gamma rays (\chref{ch:nucapps}).

\section{Averages and Probability}\label{sec:averages}

A gas contains about $10^{25}$ molecules in every cubic metre, and a milligram of a radioactive element
about $10^{18}$ nuclei. No calculation could follow each one, and none needs to: what is measured is an
average (the pressure, the temperature, the number of decays per second), and averages over huge numbers
are sharply predictable even when each individual event is not. This section collects the three tools
for working with them: averages over a distribution, the idea of a probability per unit time, and the
two integrals that the Maxwell distribution of molecular speeds requires.

\begin{defbox}{Probability, and averages over a distribution}
The \textbf{probability} $p$ of an outcome is the fraction of a very large number of trials in which it
occurs, so $0\le p\le1$, and the probabilities of all the possible outcomes add up to $1$. Two rules
combine them:
\begin{itemize}
\item for \emph{mutually exclusive} outcomes (one excludes the other), probabilities \textbf{add}: a die
      shows $1$ or $2$ with probability $\tfrac16+\tfrac16=\tfrac13$;
\item for \emph{independent} events (one does not influence the other), probabilities \textbf{multiply}:
      two dice both show $6$ with probability $\tfrac16\times\tfrac16=\tfrac1{36}$.
\end{itemize}
If a quantity $g$ takes the values $g_i$ with probabilities $p_i$, its \textbf{average} (or mean, or
expected value) is the probability-weighted sum $\langle g\rangle=\sum_i p_i\,g_i$. When the quantity
varies continuously, like a molecular speed $v$, probability is spread out instead: a
\textbf{distribution} $f(v)$ is defined so that $f(v)\,dv$ is the probability of a speed between $v$ and
$v+dv$. The sum becomes an integral,
\[
\int f(v)\,dv=1,\qquad \langle g\rangle=\int g(v)\,f(v)\,dv ,
\]
the first saying that the speed is certainly \emph{some} value (\textbf{normalisation}), the second
giving the average of any function $g$ of it. The integral is the ``sum of thin slices'' of
Section~\ref{sec:integration}, each slice being one narrow range of speeds weighted by its probability.
\end{defbox}

\begin{trapbox}{the average of a square is not the square of the average}
For the five speeds $1,2,3,4,5$ (in \si{m/s}) the mean is $\langle v\rangle=\SI{3}{m/s}$, but
$\langle v^2\rangle=(1+4+9+16+25)/5=\SI{11}{m^2/s^2}$, so the root-mean-square speed is
$\sqrt{\langle v^2\rangle}=\SI{3.32}{m/s}$, not $3$. Squaring before averaging gives the fast values extra
weight, so $\sqrt{\langle v^2\rangle}\ge\langle v\rangle$ always, with equality only if every value is the
same. Kinetic energy depends on $v^2$, which is why the kinetic theory of gases works with the
root-mean-square speed.
\end{trapbox}

The second tool concerns events that happen at random moments, such as the decay of an unstable
nucleus. What can be said about something that cannot be predicted?

\begin{derivbox}{Probability per unit time, and the survival law}
\textbf{The assumption.} An event (a decay, say) has a constant \textbf{probability per unit time}
$\lambda$: in any short interval $dt$ the chance that it happens is $\lambda\,dt$, the same at every
moment and independent of how long the system has already waited. The units of $\lambda$ are
\si{s^{-1}}. It is a rate, not a probability: only the product $\lambda\,dt$, for an interval short
enough that $\lambda\,dt\ll1$, is a probability.

\stepn{1} \textbf{Survive one more short step.} Let $S(t)$ be the probability of surviving (the event
not yet having happened) up to time $t$. To survive to $t+dt$ one must survive to $t$ \emph{and} then
survive the next $dt$, with probability $1-\lambda\,dt$. The two are independent, so the probabilities
multiply:
\[
S(t+dt)=S(t)\,(1-\lambda\,dt)\qquad\Longrightarrow\qquad \frac{dS}{dt}=-\lambda S .
\]
\stepn{2} \textbf{Solve.} This is first-order decay (Section~\ref{sec:explog}) with $S(0)=1$, so
\[
\boxed{\;S(t)=e^{-\lambda t}\;}
\]
and the probability that the event has happened by time $t$ is $1-e^{-\lambda t}$.

\stepn{3} \textbf{From one to many.} With $N_0$ independent systems, the expected number still waiting at
time $t$ is $N_0S(t)=N_0e^{-\lambda t}$. A smooth exponential law for the population emerges from random,
unpredictable individuals.

\stepn{4} \textbf{The mean waiting time.} The probability of the event falling between $t$ and $t+dt$ is
$S(t)\,\lambda\,dt$, so the average waiting time is $\langle t\rangle=\int_0^\infty t\,\lambda e^{-\lambda
t}\,dt$. The product rule confirms that
$\dfrac{d}{dt}\Big[-\Big(t+\dfrac1\lambda\Big)e^{-\lambda t}\Big]=\lambda t\,e^{-\lambda t}$, so the
integral is the bracket evaluated between the limits, $0-\big(-\tfrac1\lambda\big)$:
\[
\langle t\rangle=\frac{1}{\lambda}.
\]
The mean waiting time is the time constant of the survival curve.
\end{derivbox}

\inwords{A constant probability per unit time means the system has no memory. A nucleus that has
already survived a thousand years faces exactly the same odds in the next second as a newly formed one,
just as a fair coin that has landed heads five times is still equally likely to land tails next.
Individual outcomes stay unpredictable; only the fractions are fixed.}

\begin{exbox}{A rate is not a probability}
Each member of a large collection of independent unstable nuclei has $\lambda=\SI{0.10}{s^{-1}}$. Find
the probability that a given nucleus (a) decays within the next \SI{1.0}{s}, (b) survives the next
\SI{10}{s}, (c) decays within the next \SI{10}{s}.

\stepn{1} (a) $1-e^{-0.10}=0.095$. The shortcut $\lambda\,\Delta t=0.10$ is within $5\%$, because
$\lambda\,\Delta t$ is small.

\stepn{2} (b) $S=e^{-\lambda t}=e^{-1}=0.37$.

\stepn{3} (c) $1-e^{-1}=0.63$. Here the shortcut $\lambda\,\Delta t=1.0$ would claim certainty, which is
absurd; $\lambda\,\Delta t$ is a probability only when it is small.
\ans{$0.095$; $0.37$; $0.63$}
\end{exbox}

\textbf{Where the model stops working.} The survival law needs the events to be independent and
$\lambda$ to be truly constant; both hold extremely well for radioactive nuclei, and much less well for,
say, light bulbs, which wear out and so do age. And for a small population the counts scatter about the
average: a count whose expected value is $N$ typically differs from it by about $\sqrt N$ (stated here
without proof, because it needs the theory of the binomial distribution), a fraction $1/\sqrt N$ that is
negligible for $10^{18}$ nuclei but large for $100$.

The third tool is a pair of definite integrals. The distribution of molecular velocities in a gas,
derived in \chref{ch:kinetic}, is a bell-shaped curve proportional to $e^{-cv^2}$, where $c$ is a positive
constant, and every average over it reduces to one of the two integrals below.

\begin{keybox}{The two Gaussian integrals}
For any constant $c>0$,
\[
\boxed{\;\int_{-\infty}^{\infty}e^{-cv^2}\,dv=\sqrt{\frac{\pi}{c}},\qquad
\int_{-\infty}^{\infty}v^{2}e^{-cv^2}\,dv=\frac12\sqrt{\frac{\pi}{c^{3}}}\;}
\]
The first is quoted without proof, because its standard proof squares the integral and evaluates it as
an area integral in polar coordinates, which is beyond this book. The second follows from the first by
differentiating both sides with respect to $c$ (allowed here because the integrand is smooth and dies
away fast): the left side gives $-\int v^2e^{-cv^2}dv$, the right side gives
$\sqrt\pi\cdot\big(-\tfrac12c^{-3/2}\big)$, and cancelling the minus signs gives the result. Two
companions follow from symmetry: $\int_{-\infty}^{\infty}v\,e^{-cv^2}dv=0$, because the integrand is odd
(the areas for $v<0$ and $v>0$ cancel), and each integral from $0$ to $\infty$ of an even integrand is
half of the full one.

\textbf{The half-range moments.} Speeds are never negative, so averages over a speed distribution run from
$0$ to $\infty$. Halving the two results above, and differentiating once more with respect to $c$ for the
fourth power, gives the even moments; the odd ones need no Gaussian integral at all, because
$\frac{d}{dv}\big[-e^{-cv^2}/(2c)\big]=v\,e^{-cv^2}$ exactly, and differentiating that result with respect to
$c$ gives the next:
\[
\int_0^{\infty}e^{-cv^2}dv=\frac12\sqrt{\frac{\pi}{c}},\qquad
\int_0^{\infty}v\,e^{-cv^2}dv=\frac{1}{2c},\qquad
\int_0^{\infty}v^{2}e^{-cv^2}dv=\frac14\sqrt{\frac{\pi}{c^{3}}},
\]
\[
\int_0^{\infty}v^{3}e^{-cv^2}dv=\frac{1}{2c^{2}},\qquad
\int_0^{\infty}v^{4}e^{-cv^2}dv=\frac38\sqrt{\frac{\pi}{c^{5}}}.
\]
Each step down the list is $-d/dc$ of the entry two places above it; for example
$-\frac{d}{dc}\big(\tfrac14\sqrt\pi\,c^{-3/2}\big)=\tfrac38\sqrt\pi\,c^{-5/2}$.
\end{keybox}

\begin{exbox}{The average of $v_x^2$ for a bell-shaped distribution}
A velocity component $v_x$ is distributed as $f(v_x)=A\,e^{-cv_x^2}$. Find $A$ and $\langle v_x^2\rangle$.

\stepn{1} Normalisation: $A\sqrt{\pi/c}=1$, so $A=\sqrt{c/\pi}$.

\stepn{2} The average: $\langle v_x^2\rangle=A\int v_x^2e^{-cv_x^2}dv_x=\sqrt{\dfrac{c}{\pi}}\cdot
\dfrac12\sqrt{\dfrac{\pi}{c^{3}}}=\dfrac{1}{2c}$.

\stepn{3} In the Maxwell distribution $c=m/(2k_BT)$, where $m$ is the mass of a molecule, $T$ the
absolute temperature and $k_B$ Boltzmann's constant, so $\langle v_x^2\rangle=k_BT/m$ and
$\tfrac12m\langle v_x^2\rangle=\tfrac12k_BT$: half of $k_BT$ of kinetic energy for each direction of
motion, the result behind the equipartition of energy.
\ans{$A=\sqrt{c/\pi}$; $\langle v_x^2\rangle=1/(2c)$}
\end{exbox}

\medskip\noindent
\emph{Used in:} the mean, root-mean-square and most probable molecular speeds, the Maxwell
distribution and the equipartition of energy (\chref{ch:kinetic}); the survival law, the mean life
$\tau=1/\lambda$ and branching decays (\chref{ch:decaylaw}); and the statistics of counting decays (\chref{ch:nucapps}).

\section{Summary}\label{sec:summary0}
The chapter assembled sixteen tools, and they divide cleanly by purpose. Two \emph{describe}
directions: radian measure, the only angle unit that works inside calculus, and the trigonometric
ratios, used above all to resolve a vector. Two \emph{combine} vectors: the dot product (how much
along) and the cross product (how much across). Three connect a quantity with its rate of change: the derivative extracts an
instantaneous rate from a changing quantity (its product, chain and logarithmic forms extend it), its
zeros locate maxima and minima, and the integral adds up infinitely many thin slices (one slicing
pattern serves every body). The last five turn equations into answers: the small-angle and binomial
approximations, the quadratic formula (where both roots must be
examined for physical meaning), the elimination of an unwanted unknown between simultaneous equations,
and proportional reasoning. The final four handle change that feeds on itself and events too many to
follow one by one: the exponential and its inverse the logarithm, the first-order rate equations
$dy/dt=-ky$ and $dy/dt=P-ky$, averages over a probability distribution, and the probability per unit time
of a random event, with the two Gaussian integrals that the Maxwell distribution needs.

\begin{itemize}
\item Radians are dimensionless and $s=R\theta$; an angle standing on its own in a formula must be in
      radians.
\item $A_x=A\cos\theta$ only when $\theta$ is measured from the $x$-axis; cosine goes with the
      \emph{adjacent} leg.
\item $\vc{A}\cdot\vc{B}=AB\cos\theta$ is a number; $|\vc{A}\ex\vc{B}|=AB\sin\theta$ is the area of the
      parallelogram the two vectors span.
\item Logarithmic differentiation turns $Z=A^pB^q/C^r$ into $dZ/Z=p\,dA/A+q\,dB/B-r\,dC/C$, whose
      worst case, $\Delta Z/Z=p\,\Delta A/A+q\,\Delta B/B+r\,\Delta C/C$, is the whole theory of error
      propagation in one line.
\item A smooth quantity is greatest or least where its derivative is zero; the sign of the second
      derivative says which. For small $x$, $(1+x)^n\approx1+nx$.
\item Every integral over a body follows one pattern: choose a slice on which the weighted quantity is
      constant, express its mass, and integrate.
\item Proportional reasoning ($y\propto x^n$, so scaling $x$ by $k$ scales $y$ by $k^n$) answers many
      examination questions with no arithmetic at all.
\item $e^{x}$ is its own derivative and $\ln x$ undoes it: $\ln(ab)=\ln a+\ln b$, $d(\ln x)/dx=1/x$,
      $\int_a^b dx/x=\ln(b/a)$; $\ln2=0.693$, $\ln10=2.303$.
\item $dy/dt=-ky$ gives $y=y_0e^{-kt}$, with time constant $\tau=1/k$ and half-life
      $t_{1/2}=\ln2/k=0.693\,\tau$; $dy/dt=P-ky$ relaxes to $P/k$ with the same time constant. On a
      semilog plot an exponential is a straight line of slope $-k$.
\item Probabilities of exclusive outcomes add, of independent events multiply;
      $\langle g\rangle=\int g\,f\,dv$ with $\int f\,dv=1$, and $\langle v^2\rangle\ne\langle v\rangle^2$.
\item A constant probability per unit time $\lambda$ gives the survival probability $e^{-\lambda t}$ and
      the mean waiting time $1/\lambda$; $\lambda\,\Delta t$ is a probability only when it is small.
\item $\int_{-\infty}^{\infty}e^{-cv^2}dv=\sqrt{\pi/c}$ and
      $\int_{-\infty}^{\infty}v^2e^{-cv^2}dv=\tfrac12\sqrt{\pi/c^3}$.
\end{itemize}

\section{Problem Types}\label{sec:types0}
Mathematics is rarely examined for its own sake, but six slips recur so often inside physics questions
that they are worth isolating: an angle left in degrees, a missing chain-rule factor, a variable
quantity treated as constant, a confusion between a factor and an increase, an exponential treated as
though it fell in a straight line, and a rate mistaken for a probability. Each problem type below
gives the \emph{tell} that identifies the situation, the \emph{method}, and the \emph{trap}.

\archq{Arc length and the radian}
{An angle is given in degrees, and the formula uses it on its own: $s=R\theta$, $v=\omega r$,
$\omega=\theta/t$.}
{Convert first: $\theta_{\text{rad}}=\theta^\circ\times\pi/180$. Then substitute.}
{Substituting the number of degrees directly, which gives an answer about $57$ times too large.}
{A wheel of radius \SI{0.50}{m} turns through $60^\circ$. The distance moved by a point on its rim is:}
{\SI{30}{m}}{\SI{0.52}{m}}{\SI{0.26}{m}}{\SI{1.05}{m}\archans{B}{$s=R\theta=(0.50)(\pi/3)=\SI{0.52}{m}$.
Option (A) uses $\theta=60$ directly; option (D) uses the diameter in place of the radius.}}

\archq{Differentiating to find a rate}
{A quantity is given as a function of time, such as $x(t)$, and its rate at a particular instant is
wanted.}
{Differentiate term by term; for $\sin(\omega t)$ or $\cos(\omega t)$ the chain rule brings out a factor
$\omega$. Then substitute the instant.}
{Dividing $x$ by $t$, which gives an \emph{average} rate, or forgetting the factor $\omega$.}
{A particle moves along a line with $x=2\sin(3t)$, where $x$ is in metres and $t$ in seconds. Its
velocity at $t=0$ is:}
{$0$}{\SI{2}{m/s}}{\SI{6}{m/s}}{\SI{3}{m/s}\archans{C}{$v=dx/dt=6\cos(3t)$, which is \SI{6}{m/s} at
$t=0$. Option (A) is the position at $t=0$; option (B) omits the chain-rule factor $3$.}}

\archq{Integrating a varying quantity}
{The quantity to be summed changes along the way: the area under a curve, or the total of something
whose rate is a given function.}
{Write the sum as a definite integral and evaluate it as $G(b)-G(a)$, where $G'=f$.}
{Multiplying the final value of $f$ by the interval, as though $f$ had been constant throughout.}
{The area under the graph of $y=6x^2$ between $x=0$ and $x=2$ is:}
{$8$}{$16$}{$24$}{$48$\archans{B}{$\int_0^2 6x^2\,dx=\big[2x^3\big]_0^2=16$. Option (D) is
$y(2)\times2=24\times2$, the constant-height trap; option (C) is $y(2)$ itself.}}

\archq{Scaling by proportional reasoning}
{``By what factor (or what percentage) does $y$ change if $x$ changes by \ldots?'', with a power law
relating them.}
{If $y\propto x^n$, scaling $x$ by $k$ scales $y$ by $k^n$. Convert the answer to a percentage
\emph{increase} only at the end: an increase is $(k^n-1)\times100\%$.}
{Reporting the factor $k^n$ as though it were the percentage increase, or using $nk$ in place of $k^n$.}
{The speed of a body of fixed mass is increased by $50\%$. Its kinetic energy increases by:}
{$50\%$}{$100\%$}{$125\%$}{$225\%$\archans{C}{$K\propto v^2$ and $k=1.5$, so $K$ is multiplied by
$1.5^2=2.25$, an increase of $125\%$. Option (D) mistakes the factor for the increase; option (B) is
$n\times50\%$.}}

\archq{Exponential decay and the half-life}
{A quantity falls by the same \emph{fraction} in equal times (a half-life, a cooling time, a decay
constant is given), and the time to reach some level that is not a whole number of half-lives is wanted.}
{Write $y=y_0e^{-kt}$ with $k=\ln2/t_{1/2}$, then take natural logarithms: $t=\ln(y_0/y)/k
=t_{1/2}\,\ln(y_0/y)/\ln2$.}
{Extending the first halving as a straight line, using $t_{1/2}$ where $\tau=1/k$ belongs, or taking the
logarithm in base $10$.}
{A quantity decays exponentially with a half-life of \SI{10}{min}. The time for it to fall to $10\%$ of
its initial value is closest to:}
{\SI{33}{min}}{\SI{18}{min}}{\SI{23}{min}}{\SI{14}{min}\archans{A}{$t=t_{1/2}\ln10/\ln2
=10\times2.303/0.693=\SI{33.2}{min}$. Option (B) continues the initial rate of loss, half of $y_0$ in
\SI{10}{min}, as a straight line; option (C) treats the half-life as the time constant,
$t=t_{1/2}\ln10$; option (D) uses $\log_{10}10=1$ in place of $\ln10$.}}

\archq{Probability per unit time}
{Random, independent events (decays, collisions) with a given rate $\lambda$; the probability of an event
within, or of surviving, a stated time is wanted.}
{The survival probability is $e^{-\lambda t}$, so the probability of the event within $t$ is
$1-e^{-\lambda t}$. Use $\lambda t$ itself only when it is much less than $1$.}
{Treating $\lambda t$ as a probability when it is not small, or giving the survival probability when the
event probability is asked for.}
{Each of a large number of independent unstable nuclei has a probability per unit time of decay
$\lambda=\SI{0.20}{s^{-1}}$. The probability that a given nucleus decays within the next \SI{5.0}{s} is:}
{$1.0$}{$0.37$}{$0.50$}{$0.63$\archans{D}{$1-e^{-\lambda t}=1-e^{-1}=0.63$. Option (A) takes
$\lambda t=1$ as the probability; option (B) is the survival probability $e^{-1}$; option (C) mistakes
the mean life $1/\lambda=\SI{5}{s}$ for the half-life, which is $0.693/\lambda=\SI{3.5}{s}$.}}
